% !TeX root = ../papers/nc-p3.tex
\section{全成分のPI窓における片側の頂点拡張判定}
\label{sec:one-sided-pi-noether-criterion}
\begin{theorem}\label{thm:one-sided-pi-noether-criterion}
標準正則窓の折り畳み $B$ がPIとする。以下の条件は全て同値である。
\begin{enumerate}
\item 全ての有限次元非頂点単純 $S$ と全頂点 $k_i$ で $\operatorname{Ext}^1_B(S,k_i)=0$。
\item 全ての $S,i$ で $\operatorname{Ext}^1_B(k_i,S)=0$。
\item $B$ が右Noetherである。
\item $B$ が左Noetherである。
\item 窓は合成整域で $B$ は正規な有限生成中心上有限である。
\end{enumerate}
辺の有限性も単純の次元上界の指定も要しない。従って中心有限でないPI正則窓があれば、
ある非頂点単純から頂点へ、またある頂点から非頂点単純へ、両方の第一拡張が実在する。
この二つの非頂点単純が同じである必要はない。
\end{theorem}

\begin{proof}[一方向の消滅から各素商の有限表示まで]
第一条件を仮定する。列挙するTorでは左の頂点単純の和を $k^{(0)}$ と書く。
$\operatorname{Tor}_1^B(S,k^{(0)})$ は第一条件の双対により零となる。
任意の無限次元斉次素商 $R$ に対し、そのtrace環を $T$、可換trace環を $C$ とする。
$T$ は $C$ 上有限であり、錐の頂点以外の各閉点で $T$ の繊維の全単純は
$B$ の非頂点単純となる。ここで定理\ref{thm:pi-vertex-extension-criterion}の格子特殊化に、
一方向の消滅だけで十分である理由を記す。
非原点なら繊維の全単純が頂点のみであることは正次数traceで排除される。
もし混在すれば、一般絶対単純の安定格子の特殊繊維は両種の組成因子を持つ
（特殊化traceとCayley--Hamiltonによる同定理の証明）。
非頂点因子から頂点因子への第一Extが常に零なので、特殊繊維は頂点型だけの商と
非頂点型だけの部分加群を持つ。組成列の隣接する逆順の因子を、該当二段が分裂するたびに入れ替えればよい。
さらにこの部分加群から商への第一Extは組成列から零である。
特殊表示をこの二部分に沿って三角化し、厚み $t^n$ の逆三角成分を順に消す。
その成分は特殊部分から特殊商へのコサイクルで、良い向きの既存の非対角成分は
この逆の成分の一次式には入らない。零Extにより頂点基底を $1+t^n H$ で変えて消せる。
完備極限で一般表現に真の不変部分が残り、絶対単純性に反する。

任意の有限生成次数付き右 $T$ 加群 $N$ に対し
\[
   \dim_k\operatorname{Tor}_j^B(N,k^{(0)})<\infty\qquad (j=0,1)
\]
である。有限な左頂点Koszul複体によりこれらは有限 $C$ 加群であり、
非原点の極大イデアル $\mathfrak m$ の各厚み $N/\mathfrak m^l N$ では組成列と
第一条件によりこの両Torが零である。有限次元逆系の極限と有限 $C$ 加群の完備化の完全性で
完備化したTorも零となり、台は頂点のみである。反対側でも第0のTorだけなら
非頂点単純と頂点単純のHomの消滅だけで同じ結論を得る。
特に $T$ は $R$ の両側で有限生成である。実際第0のTorは正次数の作用を割った空間で、
次数付き中山により有限個の生成元を与える。$T$ は $R$ の有限中心化拡大で
左右Noetherなので、有限正規化拡大に対するNoether性の降下
（Formanek--Jategaonkar、定理\ref{thm:four-matrix-balanced-simples}と同じ適用）で
$R$ は左右Noetherとなる。

$R$ 自身を右 $B$ 加群として有限表示するため、右 $R$ 加群 $R$ に有限フィルトレーションを作る。
各非零残余商の中に、有限次元の非零部分加群か、ある斉次素商のtrace環上有限生成な
非零次数付き部分加群を取れる。これは定理\ref{thm:pi-vertex-extension-criterion}の
$mxT$ の抽出である。抽出に必要なのはPI商の有限素イデアル積、中心分母と、
その分母を割る際のGK次元の降下だけであり、対象が $B$ の部分加群である必要はない。
有限次元素商に着いたら一つの斉次元が生成する加群を選ぶ。
こうして部分加群を加え続ける手順は $R$ の右Noether性で止まる。
各因子について上記の第0,1のTorは有限次元である（有限次元因子には有限Koszul複体を直接使う）。
拡張でこの有限性は保たれ、特に $\operatorname{Tor}_1^B(R,k^{(0)})$ は有限次元となる。
従って $R=B/P$ の核 $P$ は右イデアルとして有限生成である。これは
$0\to P\to B\to R\to0$ のTor列と次数付き中山による。有限次元素商にもこの結論は直接成立する。
\end{proof}

\begin{proof}[根基を含む折り畳みへの復元]
有限個の斉次素イデアルの積 $P_1\cdots P_m=0$ を取る。各 $P_i$ は右有限生成で、
途中の積 $I_j=P_1\cdots P_j$ も右有限生成である。実際二つの両側イデアルの右生成元の
対の積で積イデアルを右生成できる。$I_0=B$ とすれば各 $I_{j-1}/I_j$ は
右Noether商 $B/P_j$ 上有限生成なのでNoetherである。この有限列から $B$ が右Noetherとなる。

逆に右Noetherなら、任意の非頂点単純 $S$ を整数次数 $n\ge0$ に展開した尾部
$P=\bigoplus_n S e_{n\bmod4}$ は、右 $B$ 加群として有限生成である。
その理由は、全矢を不定元 $v$ 倍した $S$ の行列表示の像に、正次数の
$c v^l 1_S$ ($c\ne0$) があることである。定理\ref{thm:four-matrix-all-regular}の軌道商の
中心多項式の証明を使う。これは像の元であって元の $B$ の中心に属する必要はない。
$l$ は4の倍数で、尾部の有限個の初期次数から生成できる。右Noether性で
この尾部の極小次数付き射影分解の各項は有限生成となる。
十分大きい終端次数で左頂点Koszul複体を尾部とテンソルしたものは、
周期表示 $S$ と各頂点の同じTor複体になる。各Torの生成次数は先ほどの有限性で有界だから
$\operatorname{Tor}_j^B(S,k^{(0)})=0$ が全ての $j$ で従う。
四次元CY双対で両方向の頂点第一Extが消え、定理\ref{thm:pi-vertex-extension-criterion}と
定理\ref{thm:noetherian-pi-window-cutoff}を適用できる。
左Noetherともう一方の消滅からの含意は反対環で同じであり、中心有限なら両側Noetherである。
\end{proof}
